\section{Propagation from small corners}
\label{sec:propagation}

Continue with the notation and hypotheses of Theorem~\ref{thm:free-corner},
and set $H=L^2(\cL)$. After Lemma~\ref{lem:free-symmetrization}, we may use
the symmetrized operator $G$ and the operators $B_k\in\mathcal B(J_k)$
from Lemma~\ref{lem:word-structure}. Our goal is to compare the action of
$B_k$ with that of $B_0$ along repeated letters, with an error independent
of $k$. We first extract a consequence of the small-corner hypothesis
that needs only the exact commutation with $L_D$.

\begin{lemma}[Full columns from a small corner]
\label{lem:full-column}
For every projection $r\in\cL$ with $\tau(r)\le t$, there is $b\in\cL$
with $b=rb$ such that
\begin{equation}
  \big\|(G-R_b)|_{Hr}\big\|\le c_0\epsilon.
  \label{eq:full-column}
\end{equation}
The same assertion holds for $G^*$.
\end{lemma}

\begin{proof}
The case $r=0$ is immediate. Fix $r\ne0$ and a projection $e\in D$ with
$\tau(e)=t$. The factoriality established in Lemma~\ref{lem:free-factor}
and comparison of projections give $r'\le e$ and a partial isometry
$v\in\cL$ such that $v^*v=r$ and $vv^*=r'$. Right multiplication by $v^*$
is a unitary left $e\cL e$-module map
\[
  eHr\longrightarrow eHr',\qquad \xi\longmapsto\xi v^*.
\]
Indeed its inverse is right multiplication by $v$, and
\[
  \|\xi v^*\|_2^2
    =\tau(v\xi^*\xi v^*)
    =\tau(\xi^*\xi r)=\|\xi\|_2^2.
\]
These equalities also hold for $L^2$ vectors by continuity. The vector $r'$
is cyclic for the left action of $e\cL e$ on $eHr'$, since
$(e\cL e)r'=e\cL r'$ is dense in that Hilbert space. Thus $eHr$ is a
cyclic domain representation of the unital algebra $e\cL e$.

Since $G$ commutes with $L_e$, its restriction defines a bounded map
$G_e:eHr\to eH$. On these domain and codomain representations the map
\[
  \Delta(a)=G_eL_a|_{eHr}-L_aG_e,\qquad a\in e\cL e,
\]
is a rectangular derivation of norm at most $\epsilon$, by
\eqref{eq:free-corner-hypothesis}. Lemma~\ref{lem:cyclic-implementation}
therefore supplies an exact intertwiner $T:eHr\to eH$ with
\[
  TL_a|_{eHr}=L_aT\quad(a\in e\cL e),\qquad
  \|G_e-T\|\le c_0\epsilon.
\]

The value $t=1/N$ allows an extension that preserves this norm. Choose
$v_1,\ldots,v_N\in D$ with
\[
  v_i^*v_i=e,\qquad v_iv_i^*\perp v_jv_j^*\ (i\ne j),\qquad
  \sum_{i=1}^N v_iv_i^*=1.
\]
The maps $(\xi_i)\mapsto\sum_i v_i\xi_i$ identify $(eHr)^N$ with $Hr$
and $(eH)^N$ with $H$ unitarily. Define
\[
  \widetilde T\left(\sum_{i=1}^N v_i\xi_i\right)
      =\sum_{i=1}^N v_iT\xi_i,\qquad \xi_i\in eHr.
\]
In these coordinates $\widetilde T$ is $T^{(N)}$, and $G|_{Hr}$ is
$G_e^{(N)}$, because $G$ commutes with each $L_{v_i}$. Consequently
\[
  \|G|_{Hr}-\widetilde T\|=\|G_e-T\|\le c_0\epsilon.
\]
For $a\in\cL$, its left action has matrix entries $L_{v_i^*av_j}$ in
both coordinate systems. Every $v_i^*av_j$ lies in $e\cL e$, so the
intertwining identity for $T$ implies
$\widetilde T L_a|_{Hr}=L_a\widetilde T$.

Extend $\widetilde T$ by zero on $H(1-r)$. Since $R_r$ commutes with all
left multiplications, this extension belongs to $L(\cL)'$ and hence is
$R_b$ for a bounded $b\in\cL$, by the tracial commutation theorem. Its
vanishing on $H(1-r)$ says $R_b=R_bR_r$, which, with the convention
$R_b\xi=\xi b$, gives $b=rb$. This proves \eqref{eq:full-column}.
Finally $G^*$ commutes with $L_D$ and
$[G^*,L_a]=-[G,L_{a^*}]^*$, so all the same hypotheses and estimates
hold for $G^*$.
\end{proof}

\begin{lemma}[Propagation along repeated letters]
\label{lem:propagation}
For every $k\ge1$ and $v\in B=L^2(A_0)$,
\begin{equation}
  \big\|B_k(s^{\otimes k}\otimes v)
       -s^{\otimes k}\otimes B_0v\big\|
       \le 2c_0\epsilon\|v\|_2.
  \label{eq:propagation}
\end{equation}
This also holds with $B_k^*$ and $B_0^*$ in place of $B_k$ and $B_0$.
\end{lemma}

\begin{proof}
Fix $k\ge1$ and first take $v\in A_0$; the zero vector is harmless.
Choose a nonzero projection $q\in D$ with $a=\tau(q)\le t$. The left
support of $qv$ is dominated by $q$. Its right support $r$ is equivalent
to its left support by polar decomposition, so $\tau(r)\le a\le t$.
Choose $b=rb$ as in Lemma~\ref{lem:full-column} for this $r$, and put
\[
  y=qvb-qB_0v.
\]
The level-zero identity in \eqref{eq:repeated-letter} gives
$G(qv)=qB_0v$, and the tensor decomposition gives
$\|qv\|_2=\sqrt a\,\|v\|_2$. It follows that
\begin{equation}
  qy=y,\qquad \|y\|_2\le c_0\epsilon\sqrt a\,\|v\|_2.
  \label{eq:propagation-error-vector}
\end{equation}
The vector $B_0v$ need only belong to $L^2(A_0)$; every expression here
uses bounded left multiplication on this vector.

Let $z$ be a Haar unitary of $qDq$ for its normalized trace, and set
$x_n=a^{-1/2}z^n$ for $n\ge1$. Then
\[
  \tau(x_n)=0,\quad \|x_n\|_2=1,\quad
  x_n^*x_n=x_nx_n^*=a^{-1}q,
  \qquad x_n,x_n^*\rightharpoonup0\ \text{in }L^2(D).
\]
The weak convergence follows, for example, from orthonormality of the
positive and negative powers. The reduced words
$\xi_n=(sx_n)^kv$ have norm $\|v\|_2$ and right support at most $r$,
because their last $x_n$ satisfies $x_nq=x_n$. With
$A_n=(sx_n)^{k-1}s$, we have
\[
  \xi_nb-(sx_n)^kB_0v=A_nx_ny.
\]
Thus \eqref{eq:full-column} yields
\begin{equation}
  \|G\xi_n-(sx_n)^kB_0v\|_2
    \le c_0\epsilon\|v\|_2+\|A_nx_ny\|_2.
  \label{eq:propagation-prefix}
\end{equation}

To estimate the second term, write \eqref{eq:word-decomposition} as
$H=E\otimes F$, where
\[
  F=\bigoplus_{l\ge0}(E^\circ)^{\otimes l}\otimes J_l,
\]
and let $P_{\mathrm{sc}}$ be projection onto $\C1\otimes F$. This projection keeps every word level
and selects the scalar component of its first $D$ factor. On this decomposition
left multiplication by $x_n$ acts on $E$. With inner products linear in
the first variable, define $\ell_n(\eta)=\langle\eta,x_n^*\rangle$.
Then
\[
  P_{\mathrm{sc}}(x_ny)=1\otimes (\ell_n\otimes I_F)y,
  \qquad \|\ell_n\|=1.
\]
For each elementary tensor this expression tends to zero in norm, since
$x_n^*\rightharpoonup0$. Approximating the fixed vector $y$ by finite
tensor sums, and using the common bound $\|\ell_n\|=1$, proves
\begin{equation}
  \|P_{\mathrm{sc}}(x_ny)\|_2\longrightarrow0.
  \label{eq:vanishing-tensor-coefficient}
\end{equation}

On $(1-P_{\mathrm{sc}})H$, multiplication by $A_n$ prepends a reduced centered word
ending in $s$ to a leading centered $D$ letter. Each prefix letter has
$L^2$ norm one, so this restriction is an isometry. On $P_{\mathrm{sc}}H$ we use
the bound $\|L_{A_n}\|\le a^{-(k-1)/2}$. Also $qy=y$ gives
\[
  \|x_ny\|_2^2
     =\langle L_{x_n^*x_n}y,y\rangle=a^{-1}\|y\|_2^2.
\]
Consequently
\[
  \|A_nx_ny\|_2
   \le \|(1-P_{\mathrm{sc}})x_ny\|_2
        +a^{-(k-1)/2}\|P_{\mathrm{sc}}x_ny\|_2,
\]
and \eqref{eq:vanishing-tensor-coefficient} implies
\[
  \limsup_n\|A_nx_ny\|_2
       \le a^{-1/2}\|y\|_2\le c_0\epsilon\|v\|_2.
\]
Here $q,k,v,b,y$ are all fixed before $n$ tends to infinity. The
possibly large prefix bound therefore disappears with the vanishing
coefficient; a convergence rate uniform in $a$ or $k$ is unnecessary.

By the full-output identity \eqref{eq:repeated-letter}, the left side
of \eqref{eq:propagation-prefix} equals
\[
  \big\|B_k(s^{\otimes k}\otimes v)
            -s^{\otimes k}\otimes B_0v\big\|
\]
for every $n$: its endpoint and repeated $D$ factors all have norm one.
Taking the limit superior proves \eqref{eq:propagation} for bounded $v$.
Density of $A_0$ in $B$ and boundedness of the operators extend it to
all $v\in B$. Applying the same argument to $G^*$, using the adjoint
identities of Lemma~\ref{lem:word-structure}, gives the final assertion.
\end{proof}

The propagation estimate now controls all word levels with the same
error. We assemble orthogonal copies of $B$ from these levels; this
will allow a single compressed commutator to test $B_0$.

Choose $e\in D$ with $\tau(e)=t$ and $p\in D$ with $e\le p$ and
$\tau(p)=1/2$; this is possible since $t\le1/2$. Set
\begin{equation}
  d=2p-1,\qquad w=ds,\qquad u=\tau_p(e)=2t,\qquad
  \tau_p=2\tau\big|_{p\cL p}.
  \label{eq:dihedral-parameters}
\end{equation}
The elements $d$ and $s$ are free centered symmetries. For $k\ge0$ define
\[
  U_k:B\longrightarrow eH,\qquad
  U_kv=t^{-1/2}ew^{-k}v=t^{-1/2}e(sd)^kv.
\]
In \eqref{eq:word-decomposition}, this is the tensor with endpoint
$t^{-1/2}e$, repeated $D$ letter $d$, and last factor
$s^{\otimes k}\otimes v$. Hence each $U_k$ is an isometry, its range
belongs to level $k$, and these ranges are mutually orthogonal. Put
\[
  \mathcal K=\overline{\bigoplus_{k\ge0}U_kB}\subset eH.
\]
Via the $U_k$, identify $\mathcal K$ with
$\ell^2(\N_0)\otimes B$, where $\N_0=\{0,1,\ldots\}$, and set
$J=I\otimes B_0$. Denote the inclusion $\mathcal K\hookrightarrow H$
by $\iota$.

\begin{lemma}[Compression of the commutator]
\label{lem:compressed-commutator}
For every $h\in e\cL e$,
\begin{equation}
  \big\|[J,\iota^*L_h\iota]\big\|
       \le (1+4c_0)\epsilon\|h\|.
  \label{eq:compression-bound}
\end{equation}
\end{lemma}

\begin{proof}
Lemma~\ref{lem:propagation} and \eqref{eq:repeated-letter} give
\[
  \|GU_kv-U_kB_0v\|_2\le2c_0\epsilon\|v\|_2,
\]
with zero error for $k=0$. Each error lies in level $k$ by
Lemma~\ref{lem:word-structure}. Their squared norms therefore add on
finite orthogonal sums. Passing to the closure, and repeating the
argument for $G^*$, gives the full-codomain estimates
\begin{equation}
  \|G\iota-\iota J\|\le2c_0\epsilon,
  \qquad \|G^*\iota-\iota J^*\|\le2c_0\epsilon.
  \label{eq:two-inclusion-errors}
\end{equation}
Write $R=G\iota-\iota J$, $R'=G^*\iota-\iota J^*$, and
$T_h=\iota^*L_h\iota$. Taking the adjoint of the second error identity
gives $\iota^*G=J\iota^*+(R')^*$. A direct expansion now yields
\[
  \iota^*[G,L_h]\iota
     =[J,T_h]+(R')^*L_h\iota-\iota^*L_hR.
\]
Apply \eqref{eq:free-corner-hypothesis} to the first term and
\eqref{eq:two-inclusion-errors} to the other two. Their bounds are
$\epsilon\|h\|$, $2c_0\epsilon\|h\|$, and
$2c_0\epsilon\|h\|$, proving \eqref{eq:compression-bound}.
\end{proof}
